\section{引言：从预训练到后训练}

本节课是 CS336 后训练（post-training）系列的第一讲，核心目标是理解如何将一个预训练模型（如 GPT-3）转化为一个可用的、安全的指令跟随模型（如 ChatGPT）。整节课的结构大致对应 InstructGPT 论文中的三步流程：监督微调（SFT）、奖励模型训练、以及基于强化学习的策略优化。

\begin{importantbox}{本节课的核心问题}
预训练赋予模型丰富的能力（推理、问答、编程等），但模型不会``开箱即用''地执行这些任务。后训练的目标是：
\begin{enumerate}
    \item 收集我们期望模型表现出的行为数据
    \item 用合适的算法让模型学会这些行为
    \item 在保证安全性的前提下规模化这一过程
\end{enumerate}
\end{importantbox}

\begin{figure}[H]
\centering
\includegraphics[width=0.95\textwidth]{slides-images/slide-001.jpg}
\caption{从 GPT-3 到 ChatGPT 的转变：左侧展示了 GPT-3 时代的应用（如 Copy.ai 生成广告文案），右侧是 ChatGPT 带来的对话革命\protect\footnotemark}
\end{figure}
\footnotetext{Slides 第2页。}

GPT-3 虽然参数量巨大、能力丰富，但在产品层面并不实用——它不遵循指令、不进行对话、也不拒绝恶意请求。ChatGPT 的出现改变了这一切，其核心在于后训练流程的成功应用。

\begin{figure}[H]
\centering
\includegraphics[width=0.95\textwidth]{slides-images/slide-002.jpg}
\caption{指令跟随能力的惊人之处：GPT-4 能够零样本地遵循复杂的嵌套复合指令，生成 matplotlib 代码并绘制精确图表\protect\footnotemark}
\end{figure}
\footnotetext{Slides 第3页。示例来自 Bubeck et al. 2023, ``Sparks of AGI''.}

\begin{figure}[H]
\centering
\includegraphics[width=0.95\textwidth]{slides-images/slide-003.jpg}
\caption{安全与内容审核的必要性：左图展示输入过滤器拦截仇恨言论请求；右图展示更复杂的社会工程攻击中输入和输出过滤器的协同工作\protect\footnotemark}
\end{figure}
\footnotetext{Slides 第4页。}

\begin{knowledgebox}{预训练 vs. 后训练的心智模型}
可以这样理解两者的关系：
\begin{itemize}
    \item \textbf{预训练}：将各种能力``装载''到模型参数中——模型在某处``知道''如何推理、回答问题等
    \item \textbf{后训练}：让模型将这些能力``开箱即用''地展现出来——通过少量高质量数据引导行为
\end{itemize}
这一类比在后续讨论``SFT 能否教授新知识''和``幻觉问题''时会被反复提及。
\end{knowledgebox}

\begin{figure}[H]
\centering
\includegraphics[width=0.95\textwidth]{slides-images/slide-005.jpg}
\caption{InstructGPT 的三步流程：Step 1 收集示范数据做 SFT；Step 2 收集比较数据训练奖励模型；Step 3 用 PPO 优化策略\protect\footnotemark}
\end{figure}
\footnotetext{Slides 第6页。来自 Ouyang et al. 2022.}


